% sec-00: front matter.  Abstract and the reading path.
\section*{Abstract}
\addcontentsline{toc}{section}{Abstract}

In the week of September 28, 2026 ChemicalQDevice answered four signals in four
days: two unsolicited LinkedIn inquiries from job applicants, the SEC San
Francisco Regional Office's word that it is taking the chief executive's request
``very seriously'', several White House communications after the September 12
message, and a National Science Foundation inquiry about applying to a
multi-million-dollar pilot program. Sixteen letters, two LinkedIn replies and
three submissions went out. This packet is the fifth day, and it spends no new
signal.

It records every letter of the week against one rule. A letter earns one
follow-up on the third business day after it was sent, counting only days the
recipient's office was open, and only if no reply has arrived and the follow-up
carries one new fact. By that rule, four letters asked for nothing and earn
nothing, six have not yet earned a follow-up, two share one, one is held and one
is superseded, and three follow-ups leave today: one to the National Cancer
Institute's SBIR Development Center with a one-page Specific Aims
\cite{ncisbir2026,ncisbirpd2026}, and one each to the SEC office
\cite{secsf2026} and the FDA's Oncology AI Program \cite{fdaoceai2026}. The
packet then states the robotic Phase 1 behind every letter in numbers a funder
can check, and lists the funding routes the week opened, with the amount and the
next step on each.

\subsection*{How to read this}

A reader with five minutes takes \S1, the week as one record, and \S2, the day's
actions and the weekly cadence the day adopts.

A reader asking why only three letters go today takes \S3, the follow-up rule and
the ledger of all sixteen letters.

A funder takes \S4, the robotic Phase 1 in numbers, and \S5, the routes opened
this week. \S6 carries the positioning constraints, the method, and the
references, every one of which is a clickable link.
